% from ega4-2-fr.tex
\oldpage[IV-2]{197}

extension finie $k$ de $k(x)$ et tout point $z$ de
$f^{-1}(x)\otimes_{k(x)}k=Z$, la propriété
$\boldsymbol R(\mathscr O_{Z,z})$ soit vraie, il faut et il suffit que la
condition suivante soit satisfaite : pour toute $A$-algèbre intègre finie
$B$, si
$g:\operatorname{Spec}(A'\otimes_A B)\to\operatorname{Spec}(B)$ est le
morphisme déduit de $f$ par extension à $B$ de l'anneau de base, et si
$T=g^{-1}(\xi)$ est la fibre du point générique $\xi$ de
$\operatorname{Spec}(B)$, alors, pour tout $t\in T$,
$\boldsymbol R(\mathscr O_{T,t})$ est vraie.
\end{lemma}

\begin{proof}
La condition est trivialement nécessaire puisque si $x$ est le point de
$\operatorname{Spec}(A)$ au-dessus duquel est $\xi$, $k(\xi)$ est une
extension finie de $k(x)$. Inversement, considérons un point quelconque
$x\in\operatorname{Spec}(A)$ et soit $k$ une extension finie de $k(x)$.
Si l'on pose $\mathfrak p=\mathfrak J_x$, $k(x)$ est le corps des
fractions de $A/\mathfrak p$, et il y a une base de $k$ sur $k(x)$ formée
d'éléments entiers sur $A/\mathfrak p$; si $B$ est le sous-anneau de $k$
engendré par ces éléments, $B$ est une $A$-algèbre intègre finie et $k$
est le corps $k(\xi)$ au point générique $\xi$ de
$\operatorname{Spec}(B)$; comme $x$ est l'image de $\xi$ dans
$\operatorname{Spec}(A)$, la fibre $g^{-1}(\xi)$ n'est autre que
$f^{-1}(x)\otimes_{k(x)}k$, ce qui achève de démontrer le lemme.
\end{proof}

\begin{proof}[Fin de la démonstration de (7.3.16)]
Le fait que a) entraîne c) est une conséquence immédiate de (7.3.15) et
(7.3.12). D'autre part, si on spécialise c) au cas où $B$ est intègre, et
si $y$ désigne le point générique de $Y=\operatorname{Spec}(B)$, on a
$\mathscr O_{Y,y}=k(y)$, donc $y\in U_{\boldsymbol R}(Y)$, puisque
$\boldsymbol R(k)$ est vraie pour tout corps $k$; en exprimant que tout
point $y'\in g^{-1}(y)$ appartient à $U_{\boldsymbol R}(Y')$, on obtient
l'énoncé b), donc c) entraîne b). C.Q.F.D.
\end{proof}

\begin{proposition}[7.3.17]
\label{IV.2.s7.3.17-fr}
Supposons que la propriété $\boldsymbol R$ vérifie les conditions
$(\mathrm R'_{\mathrm I})$ et $(\mathrm R_{\mathrm{II}})$ et que
$\boldsymbol P$ soit la propriété du premier (resp. du second) type
définie à partir de $\boldsymbol R$ (7.3.10). Alors, si $A$ est un
$\boldsymbol P$-anneau, les propriétés $\boldsymbol R(A)$ et
$\boldsymbol R(\widehat A)$ sont équivalentes.
\end{proposition}

\begin{proof}
Cela résulte de (7.3.12.1) appliqué à
$Y=\operatorname{Spec}(A)$ et $X=\operatorname{Spec}(\widehat A)$.
\end{proof}

\begin{proposition}[7.3.18]
\label{IV.2.s7.3.18-fr}
Supposons que $\boldsymbol P$ soit une propriété du premier (resp. du
second) type définie à partir d'une propriété $\boldsymbol R$ vérifiant
les conditions $(\mathrm R'_{\mathrm I})$ et
$(\mathrm R_{\mathrm{II}})$ ainsi que la condition suivante :

$(\mathrm R_{\mathrm{III}})$ Pour tout anneau local noethérien complet
$C$, si l'on pose $Z=\operatorname{Spec}(C)$, l'ensemble
$U_{\boldsymbol R}(Z)$ (7.3.12) est ouvert dans $Z$.

Alors, si $A$ est un $\boldsymbol P$-anneau et si l'on pose
$X=\operatorname{Spec}(A)$, l'ensemble $U_{\boldsymbol R}(X)$ est ouvert
dans $X$.
\end{proposition}

\begin{proof}
En effet, si $X'=\operatorname{Spec}(\widehat A)$ et si $f:X'\to X$ est
le morphisme canonique, on a
$U_{\boldsymbol R}(X')=f^{-1}(U_{\boldsymbol R}(X))$ (7.3.12). Comme $f$
est fidèlement plat et quasi-compact et que par hypothèse
$U_{\boldsymbol R}(X')$ est ouvert dans $X'$, la conclusion découle de
(2.3.12).
\end{proof}

\begin{remarks}[7.3.19]
\EGAReaderTarget{ega4.c04.statement.7-3-19}{IV.2.s7.3.19-fr}
\begin{enumerate}
\item[\emph{i)}] La propriété $\boldsymbol R$ qui sert à définir $\boldsymbol P$
vérifie la condition $(\mathrm R_{\mathrm{III}})$ dans tous les exemples
énumérés dans (7.3.8). Pour (i), (ii) et (iii), cela résulte de (6.11.2)
et du théorème de Cohen (0, 19.8.8); lorsque $\boldsymbol P$ est une des
propriétés (iv), (iv$'$), (v), (v$'$), cela découle de (6.12.7) et
(6.12.9), et lorsque $\boldsymbol P$ est l'une des propriétés (vii),
(vii$'$), de (6.12.7) et (6.13.4); enfin, pour (vi) et (vi$'$)
l'assertion de (7.3.18) est triviale, étant vraie pour tout préschéma
localement noethérien.

